% !TEX program = lualatex
% Fictional sample inspired by Star Wars. Keep personal data in private/.
\documentclass[a4paper]{ltjsarticle}
\usepackage{career}

\newcommand{\CareerDate}{\number\year 年\number\month 月\number\day 日}
\newcommand{\CareerName}{ヨーダ}

\begin{document}
\CareerHeader{\CareerDate}{\CareerName}

\CareerSection{経歴概要}
ジェダイ騎士団で若手の育成、紛争の調停、銀河の平和維持に従事。ジェダイ・マスターとして、さまざまな立場の関係者と対話し、困難な局面における判断とチームの支援を担っています。以下は作品に着想を得た架空のサンプルです。

\CareerSection{活かせる経験・知識・技術}
\CareerBullet{パダワンの育成と、長期的な視点に立った指導}
\CareerBullet{対話を通じた紛争の調停と関係者間の合意形成}
\CareerBullet{フォースを用いた状況把握と冷静な意思決定}
\CareerBullet{ジェダイ評議会における方針検討とチーム支援}

\CareerSection{職務経歴}
\CareerCompany{ジェダイ騎士団（架空）}{1980年4月～現在}{ジェダイ評議会}

\begin{CareerProject}{2024年4月\par ～\par 現在}
  {銀河の平和維持と紛争調停}
  {\CareerField{人数}{6名}
   \CareerField{役割}{ジェダイ・マスター}
   \CareerField{利用技術・道具}{フォース\par ライトセーバー}}
\CareerHeading{【プロジェクト概要】}
銀河各地の関係者と協力し、対話による問題解決と平和維持に取り組む。
\CareerHeading{【担当業務】}
\CareerBullet{関係者の状況と要望を整理し、調停方針を検討}
\CareerBullet{評議会と連携した任務の計画・実施}
\CareerBullet{若手ジェダイへの助言と現場での支援}
\CareerHeading{【成果・貢献】}
\CareerBullet{異なる立場の意見を丁寧に聞き、対話による解決を支援。}
\CareerBullet{経験を共有し、任務を通じた若手の成長に貢献。}
\end{CareerProject}

\begin{CareerProject}{2022年4月\par ～\par 2024年3月}
  {パダワン向け修練の企画・指導}
  {\CareerField{人数}{3名}
   \CareerField{役割}{指導者}
   \CareerField{利用技術}{フォース\par 瞑想}}
\CareerHeading{【プロジェクト概要】}
若手ジェダイの基礎力と判断力を育む修練を企画し、指導を担当。
\CareerHeading{【担当業務】}
\CareerBullet{修練計画と実践課題の作成}
\CareerBullet{瞑想や対話を通じた個別指導}
\CareerHeading{【成果・貢献】}
\CareerBullet{一人ひとりの成長に合わせた指導と経験の共有を推進。}
\end{CareerProject}

\CareerSection{自己PR}
\CareerHeading{課題整理と合意形成}
立場の異なる相手の声に耳を傾け、冷静に状況を見極めながら合意を目指します。

\CareerHeading{継続的な学習}
学び続ける姿勢を忘れず、経験を次世代へ伝えることを大切にしています。
\par\hfill 以上
\end{document}
